\documentclass[10pt,journal]{IEEEtran}

% ── Packages ─────────────────────────────────────────────────────────────────
\usepackage{cite}
\usepackage{amsmath,amssymb,amsfonts}
\usepackage{graphicx}
\usepackage{booktabs}
\usepackage{array}
\usepackage{microtype}
\usepackage{url}
\usepackage[T1]{fontenc}
\usepackage{xcolor}
\usepackage{balance}

% ── Metadata ──────────────────────────────────────────────────────────────────
\title{Generative Modelling for Realistic Mars Landing Site Visualisation:\\
       A Critical Survey of the Enabling Literature}

\author{%
  \IEEEauthorblockN{Deepakraj Rajmohan}\\
  \IEEEauthorblockA{\textit{Msc in Artificial Intelligence}\\
  University of Surrey\\
  deepak.dr036@gmail.com}
}

\markboth{Literature Review}{}

% =============================================================================
\begin{document}
\maketitle

% ── Abstract ─────────────────────────────────────────────────────────────────
\begin{abstract}
Preparing rover operators for planetary surface missions demands photorealistic
virtual training environments that replicate the target landing site at the
resolution rover cameras actually deliver. For Mars, the highest-resolution
orbital resource is the High Resolution Imaging Science Experiment (HiRISE) at
25\,cm/pixel in nadir view, whereas rover cameras resolve sub-millimetre surface
texture from near-grazing perspective angles, a gap of roughly three orders of
magnitude that standard single-image super-resolution methods, which upscale at
most $4\times$, cannot close. This review critically examines twenty publications
spanning single-image super-resolution, image-to-image domain translation,
digital terrain model generation, neural radiance field rendering, and virtual
reality reconstruction, evaluating each against the specific challenge of
synthesising rover-perspective surface textures from nadir orbital imagery for
mission training simulation. The analysis shows that while each sub-field has
advanced substantially since 2021 transformer-based and diffusion-based
super-resolution outperform earlier GAN methods; NeRF representations of Mars
surface have been demonstrated from rover imagery; and physics-constrained
generative networks reduce terrain reconstruction error by more than
50\% no published method jointly addresses the view-angle transformation and
high-fidelity texture synthesis within a single framework. The three primary
barriers are the absence of aligned cross-view training data, the documented
risk of hallucination in diffusion-model outputs for mission-critical
visualisations, and the computational cost of diffusion-based synthesis at
orbital-to-rover scale. These findings delineate the specific research gap that
motivates the proposed dissertation.
\end{abstract}

\begin{IEEEkeywords}
Mars surface imagery, super-resolution, image-to-image translation, neural
radiance fields, digital terrain models, diffusion models, virtual reality
simulation, HiRISE, domain adaptation, planetary exploration
\end{IEEEkeywords}

% =============================================================================
\section{Introduction}
\label{sec:intro}

Successful deployment of a planetary rover begins not with launch, but with
exhaustive ground simulation. The ESA ExoMars Rosalind Franklin Mission (RFM),
targeting Oxia Planum as its landing site, requires training environments
sufficiently faithful to the actual terrain to rehearse traverse planning,
instrument operations, and hazard response. Virtual reality (VR) satisfies this
requirement in principle, but its operational value collapses when surface
textures become visually implausible. A crew navigating a smoothed, geometrically
approximate mesh is not rehearsing the same perceptual and decision-making
conditions it will face on the ground. The environment must look like Mars at
the scale rover cameras can actually resolve.

The core difficulty is a data paradox embedded in planetary imaging architecture.
The highest-resolution orbital survey instrument on Mars is the High Resolution
Imaging Science Experiment (HiRISE) aboard NASA's Mars Reconnaissance Orbiter,
which acquires imagery at 25--50\,cm/pixel in a nadir (directly downward)
viewing geometry~\cite{tao2021madnet}. The ExoMars Trace Gas Orbiter's Colour
and Stereo Surface Imaging System (CaSSIS) provides broader multi-spectral
coverage at a coarser 4--5\,m/pixel~\cite{tao2021cassis}. The
Mastcam-Z (Rover camera) aboard NASA's Perseverance rover being the most capable current
example resolve individual regolith grains at sub-millimetre scale from a
vantage point roughly 2\,m above the surface, looking across it at oblique
angles. The gap between HiRISE's 25\,cm/pixel and sub-millimetre rover
resolution spans approximately three orders of magnitude in linear scale, or
six orders of magnitude in area. Current single-image super-resolution (SISR)
methods reduce this gap by at most a factor of four~\cite{tao2021cassis},
leaving it essentially intact.

Resolution alone, however, does not capture the full dimensionality of the
problem. A nadir HiRISE image of a boulder field depicts the horizontal top
surfaces of rocks. A rover standing 10\,m away sees the same boulders in
perspective: vertical faces, shadows cast across the regolith floor, and the
parallax between foreground grains and background outcrops. Upsampling a
nadir image, however aggressively, cannot hallucinate these perspective-view
features because they are geometrically absent from the input. The view-angle
transformation from nadir to perspective requires inferring three-dimensional
surface geometry from two-dimensional orbital reflectance data and re-rendering
that geometry from a different camera position is a problem that SR, by
construction, does not address.

This review surveys twenty publications across seven thematic areas, assessing
each against the specific research requirement: generating rover-perspective
surface texture from nadir orbital imagery at VR-grade fidelity for pre-mission
simulation. The seven thematic areas are:
\begin{enumerate}
  \item Orbital imaging systems and Mars data characterisation;
  \item Single-image super-resolution for planetary remote sensing;
  \item Image-to-image translation frameworks applicable to view-domain bridging;
  \item Monocular digital terrain model (DTM) generation from orbital imagery;
  \item Neural radiance field (NeRF) and Gaussian splatting methods applied to
        Mars;
  \item Virtual reality reconstruction pipelines for planetary mission
        simulation;
  \item Hallucination and epistemic reliability in generative models used for
        scientific inference.
\end{enumerate}
The critical finding, developed throughout the analysis, is that no published
method and no obvious composition of published methods, closes the
nadir-orbital to rover-perspective gap within a single trainable framework.
The dissertation work targets precisely this gap.

% =============================================================================
\section{Taxonomy and Conceptual Framework}
\label{sec:taxonomy}

The literature reviewed here does not constitute a unified research programme.
It comprises several independently advancing communities like remote sensing image
processing, computer vision generative modelling, planetary science, and VR
engineering whose outputs are jointly necessary to the downstream VR synthesis
requirement without any of them addressing it explicitly. Three axes organise the
territory.

The first axis is \textit{spatial view}: methods operate in the nadir domain
(orbital, top-down), the perspective domain (rover, oblique), or across both.
The second axis is \textit{task type}: methods either enhance existing imagery
(super-resolution, inpainting), translate between domains (image-to-image
translation, domain adaptation), or recover and render geometric structure (DTM
generation, NeRF, Gaussian splatting). The third axis is \textit{output
modality}: visual RGB texture, elevation/depth maps, or rendered video sequences.

Seven clusters emerge from this structure. \textbf{Cluster~1} (Orbital data
characterisation) establishes baseline understanding of Mars instrument
capabilities and their information limits~\cite{serrano2013hirise,
osadnik2025mcted}. \textbf{Cluster~2} (Nadir SR) comprises methods that upscale
orbital images while remaining in the nadir domain~\cite{tao2021cassis,
wu2024rstsrn, lu2025thermal, khanna2024diffusionsat, zhu2025taming}.
\textbf{Cluster~3} (Image-to-image translation) contains domain-agnostic
frameworks that map between arbitrary distributions and could, in principle,
translate nadir to perspective~\cite{kim2024unsb, xiao2025brownian}.
\textbf{Cluster~4} (Monocular DTM) covers methods that infer 3D topography from
single orbital images~\cite{tao2021madnet, zou2026pfmgan, cao2024gen3d}.
\textbf{Cluster~5} (Neural rendering) encompasses NeRF and Gaussian splatting
methods applied to Mars scenes~\cite{giusti2022marf, li2025martian,
dealmeida2025neural3d, martinezpetersen2025hifi, paul2025nerfspace}.
\textbf{Cluster~6} (VR reconstruction) covers end-to-end pipelines targeting
mission simulation environments~\cite{catalano2025inpainting, li2025martian}.
\textbf{Cluster~7} (Reliability) analyses hallucination and epistemic validity
in generative scientific models~\cite{aithal2024hallucination,
rathkopf2025reliability}.

The gap between Cluster~2 and Cluster~5 is the primary concern of this
dissertation. Nadir SR methods keep the viewpoint fixed; neural rendering
methods require perspective imagery as input. No reviewed method translates from
nadir to perspective in a way that synthesises rover-viewable texture from
orbital data. This structural gap is not an oversight in the literature; it
reflects the distinct communities that have developed these tools and the absence
of a formal problem statement that spans both.

% =============================================================================
\section{Critical Analysis of the Literature}
\label{sec:analysis}

\subsection{The Orbital Data Landscape and Its Inherent Constraints}
\label{subsec:data}

HiRISE provides the sharpest available view of the Martian surface from orbit,
yet its coverage is drastically limited. Fewer than 3\% of the Martian surface
has been imaged at 25\,cm/pixel resolution, because the instrument's 6\,km swath
width and the competing demands of the science team constrain acquisition
rates~\cite{tao2021madnet}. CaSSIS covers broader areas at 4--5\,m/pixel, and
the Context Camera (CTX) on MRO provides near-global coverage at 6\,m/pixel.
Each instrument occupies a distinct rung in the resolution hierarchy, and the
gulf between even the finest orbital product and rover cameras remains
unbridged regardless of which instrument is the starting point.

Serrano \textit{et al.}~\cite{serrano2013hirise} expose a revealing consequence
of this hierarchy. The paper trains a Bayesian network to infer HiRISE-equivalent
rock density, a key landing hazard metric from coarser CTX image features,
incorporating geomorphic unit labels as prior information. The model succeeds in
probabilistic density estimation but does so without reconstructing surface
\textit{appearance}: it predicts rock abundance, not rock texture or shape. This
is not a design limitation; it is a statement about information content. CTX
pixels do not contain the sub-pixel spectral and spatial variation from which
HiRISE-equivalent visual texture could be deterministically reconstructed. Any
generative model that attempts to synthesise rover-scale visual texture from
nadir orbital imagery faces the same information-theoretic constraint. The
synthesis is necessarily generative and grounded in statistical priors rather
than purely reconstructive.

Osadnik \textit{et al.}~\cite{osadnik2025mcted} respond to data scarcity by
releasing MCTED, comprising 80,898 paired CTX orthoimage--digital elevation
model (DEM) samples derived from the Day \textit{et al.} processing pipeline,
with tooling to handle common artefacts and missing-data regions. The contribution
is genuine: standardised, ML-ready paired data has been conspicuously absent from
Mars terrain modelling research, and MCTED provides a reproducible training
benchmark. Its limitation, from the perspective of this dissertation, is
structural: the pairs link nadir images to elevation maps, not nadir images to
perspective views. A dataset providing the same surface patch in both nadir
orbital and rover perspective at aligned resolution does not yet exist, and its
absence constrains every approach in Clusters~3 and 5 as much as it constrains
novel methods.

\subsection{Super-Resolution for Planetary Remote Sensing}
\label{subsec:sr}

SR research on planetary imagery has progressed through three generations since
2021: GAN-based methods modelled closely on terrestrial benchmarks, attention and
transformer architectures that better capture long-range spatial dependencies, and
diffusion-based synthesis that trades pixel-level reconstruction accuracy for
perceptual realism.

\subsubsection{GAN-Based SR}

Tao \textit{et al.}~\cite{tao2021cassis} propose MARSGAN, an enhancement of
ESRGAN incorporating adaptive-weighted residual dense blocks (AW-RRDB) with
Gaussian noise injection and a multi-scale reconstruction scheme that uses
pixel-shuffling to suppress checkerboard artefacts. Trained exclusively on
HiRISE images and applied at inference to CaSSIS tiles, MARSGAN achieves
approximately $3\times$ effective resolution improvement, reducing apparent CaSSIS
resolution from 4--5\,m/pixel to roughly 1.3\,m/pixel. The work is evaluated
across eight geologically diverse Mars science sites including the Mars 2020
Perseverance landing site at Jezero Crater. Notably, however, the evaluation is
entirely qualitative, no PSNR or SSIM scores are reported. This is not merely a
presentation choice. Co-registering CaSSIS and HiRISE acquisitions of the same
location precisely enough for pixel-level comparison is non-trivial given the
instruments' different viewing geometries, acquisition times, and spatial
sampling. The qualitative approach is honest but leaves cross-method performance
comparison opaque, a problem that recurs throughout this sub-field.

\subsubsection{Transformer-Based SR}

Wu \textit{et al.}~\cite{wu2024rstsrn} address SR at $\times$2, $\times$4, and
$\times$8 magnification through RSTSRN, which extends LapSRN's Laplacian
pyramid architecture with recursive Swin Transformer attention blocks.
The recursive strategy processes a single input through shared attention weights
at multiple scales, reducing parameter count while expanding the effective
receptive field. On Mars image benchmarks, RSTSRN improves upon LapSRN by
0.35\,dB PSNR at $\times$2, 0.88\,dB at $\times$4, and 1.22\,dB at $\times$8,
with corresponding SSIM gains of 0.0048, 0.0114, and 0.0149, and outperforms
competing transformer architectures including SwinIR. The improvement at higher
scale factors is the most operationally relevant finding: the larger the SR
factor, the more the texture generation task resembles free synthesis rather than
reconstruction, and transformer self-attention mechanisms are better equipped
than convolutional filters to maintain long-range consistency across
synthesised regions. The critical flaw is that evaluation uses bicubic
downsampling to generate low-resolution inputs from the original images, then
measures reconstruction against those originals. This protocol assesses recovery
from one specific degradation model. Real CaSSIS or CTX pixels do not
approximate bicubically downsampled HiRISE, they emerge from a physically
distinct optical system with its own point spread function, noise floor, and
spectral response. The trained model has never encountered the actual low-to-high
resolution relationship it is meant to solve.

\subsubsection{Diffusion-Based SR and Generative Foundation Models}

Khanna \textit{et al.}~\cite{khanna2024diffusionsat} step back from
Mars-specific engineering to propose DiffusionSat, a generative foundation model
for satellite imagery broadly. Trained on a heterogeneous collection of public
remote sensing datasets including NAIP (1\,m/pixel aerial imagery of the USA),
xView, FMOW, and SpaceNet, DiffusionSat conditions generation on geolocation
metadata and sensor parameters through a 3D positional encoding scheme applied to
the diffusion model's conditioning stack. The model achieves competitive FID scores
on multi-sensor satellite image generation tasks and demonstrates, crucially, that
a single diffusion model can generalise across diverse Earth observation sensors
without sensor-specific fine-tuning. The indirect relevance to Mars is significant:
DiffusionSat shows that multi-sensor domain generalisation is tractable within a
diffusion framework. Its failure mode for Mars is equally clear: its training
corpus contains no Mars data; its geolocation conditioning assumes Earth
coordinates; and its multi-spectral handling assumes terrestrial atmospheric
spectral profiles. Transferring DiffusionSat to Mars would require either
large-scale domain fine-tuning or an architectural reformulation of the metadata
conditioning that replaces Earth-specific priors with Martian photometric models.

Zhu \textit{et al.}~\cite{zhu2025taming} address a tension that DiffusionSat
sidesteps: unconditional diffusion models generate perceptually sharp outputs but
diverge from pixel-level ground truth, while deterministic regression models
achieve high PSNR at the cost of blurred high-frequency texture. Their approach
introduces a frequency-decomposition conditioning mechanism that anchors the
diffusion denoising trajectory to low-frequency structural content from the input,
allowing high-frequency texture to be generated stochastically without
compromising large-scale geometry. Evaluated on the DOTA and DIOR-R aerial
imagery benchmarks at $\times$4 SR, the method improves PSNR over prior
diffusion SR methods while maintaining competitive FID scores. For Mars
applications, the architectural insight translates directly: any synthesis pipeline
that must preserve large-scale geomorphology (crater rims, dune crests, valley
networks) while generating sub-metre surface texture will encounter exactly
this fidelity-perceptual realism trade-off, and frequency decomposition offers
a principled way to navigate it.

Lu and Su~\cite{lu2025thermal} take a multi-source fusion approach to SR on Mars
thermal infrared imagery from the THEMIS instrument. By fusing visible-channel
orbital images and MOLA topographic data to guide thermal channel upsampling,
the method leverages the complementary spatial information carried by each
modality. The approach demonstrates measurable PSNR gains on THEMIS data.
Its relevance to the dissertation is indirect: thermal imagery is not the visual
RGB texture domain required for photorealistic VR, and the multi-source fusion
strategy while promising requires co-registered multi-modal data at the
inference site, which may not be uniformly available across candidate Mars landing
areas.

Collectively, the SR literature demonstrates that $\times$4 resolution enhancement
of orbital Mars imagery is achievable with GAN, transformer, or diffusion
backbones, and that diffusion-based methods best handle the texture stochasticity
of geological surfaces. However, $\times$4 closes a factor-of-four fraction of a
three-order-of-magnitude gap. More fundamentally, every reviewed SR method
operates nadir-to-nadir: the input is an orbital image; the output is a
higher-resolution orbital image. The perspective transformation that converts
nadir imagery into what a rover would see is nowhere addressed.

\subsection{Image-to-Image Translation as a Domain-Bridging Mechanism}
\label{subsec:i2i}

Image-to-image translation methods map entire image domains to each other, a
capability structurally better suited to the nadir-to-perspective view
transformation than SR, which merely sharpens detail within a fixed viewpoint.
Two principled frameworks appear in the reviewed literature, both grounded in
stochastic differential equation (SDE) dynamics.

Kim \textit{et al.}~\cite{kim2024unsb} propose UNSB (Unpaired Neural
Schr{\"o}dinger Bridge), which reframes unpaired domain translation as learning
an SDE that transports samples from one arbitrary distribution to another. Unlike
standard diffusion models, which route data through a Gaussian noise prior and
back, the Schr{\"o}dinger Bridge transports directly between two target
distributions. This eliminates the Gaussian prior assumption that forces standard
diffusion I2I methods to destroy domain-specific structure before rebuilding it.
UNSB addresses the scalability of the Schr{\"o}dinger Bridge formulation by
decomposing the SDE into a sequence of adversarial sub-problems, allowing
high-resolution unpaired I2I to succeed where prior SB formulations stalled.
On standard benchmarks including horse-to-zebra, cat-to-dog, summer-to-winter,
and day-to-night translations, UNSB achieves lower FID than CycleGAN, CUT, and
related methods. The architecture is compelling for the Mars nadir-to-perspective
problem precisely because it operates on unpaired data: a source domain of
HiRISE nadir patches and a target domain of rover perspective images are both
available, and they do not need to be geographically co-registered. However,
UNSB has been evaluated only on natural image colour-domain translations. The
photometric differences between a nadir view and a perspective view of the
same Mars surface far exceed the hue and seasonal saturation shifts of the tested
benchmarks; they involve changes in shadow geometry, rock face visibility,
foreshortening, and inter-object occlusion. All of which require implicit 3D
reasoning that the SDE formulation does not explicitly encode.

Xiao \textit{et al.}~\cite{xiao2025brownian} address the complementary problem:
stochastic I2I methods produce different outputs from the same input on
successive runs, which is unacceptable in scientific visualisation. Their
Dual-approx Bridge exploits Brownian bridge dynamics with two separate neural
approximators, one for the forward noising process, one for the reverse to
achieve deterministic I2I: the same input produces the same output every time.
On paired image generation and super-resolution benchmarks, Dual-approx Bridge
achieves superior SSIM and LPIPS compared to stochastic diffusion I2I baselines
while maintaining image fidelity comparable to the ground truth. Determinism
matters greatly in the Mars context: if two planetary scientists run the same
generation pipeline and obtain different reconstructions of the same surface
patch, neither output can be used as mission-planning ground truth. Dual-approx
Bridge provides the architectural template for a reproducible synthesis system,
though, like UNSB, it has not been applied to remote sensing or view-angle
translation tasks.

The gap here is categorical rather than incremental. No I2I method has been
applied to the nadir-to-perspective Mars translation problem. Neither framework
has been tested on remote sensing imagery. The photometric complexity of the
view-angle transformation is qualitatively different from the colour and seasonal
domain shifts these methods were designed and benchmarked on. Whether SB-based
I2I is expressive enough to capture geometry-dependent viewpoint changes or
whether a geometry-aware architecture is necessary remains an open empirical
question.

\subsection{Digital Terrain Model Generation from Monocular Orbital Imagery}
\label{subsec:dtm}

DTM generation from orbital imagery occupies a structurally intermediate
position: it extracts 3D geometric information from nadir images, which is a
necessary precursor to perspective-view synthesis but does not itself complete
the journey to VR-ready visual texture.

Tao \textit{et al.}~\cite{tao2021madnet} present MADNet~2.0, which retrieves
pixel-scale topography from a single HiRISE image by exploiting photoclinometric
cues, the relationship between pixel intensity, surface slope, and sun angle
at acquisition. A deep learning model trained on paired HiRISE image--MOLA
altimetry data infers a dense elevation map at HiRISE's native 25\,cm/pixel
resolution. This monocularity is the key engineering advantage: stereo-based
DTM generation requires two overlapping acquisitions of the same location, which
covers less than 1\% of the HiRISE archive; MADNet~2.0 operates on any
single-strip observation. The limitation for VR is categorical: the DTM provides
surface geometry, not surface appearance. Rocks with identical elevation profiles
but different mineralogy, grain size, or weathering state produce entirely
different visual textures, none of which MADNet~2.0 recovers.

Zou \textit{et al.}~\cite{zou2026pfmgan} address the photometric ambiguity
that confounds DTM generation through PFMGAN, which integrates an explicit solar
geometry model and an albedo-aware attention module into a GAN architecture. By
modelling how sun angle and surface albedo interact with topography to produce
the observed shading, PFMGAN reduces the confusion between shading caused by
slope and shading caused by reflectance variation. Evaluated on five Martian
landform classes (impact craters, aeolian dunes, lava flow fields, valley
networks, and plains), PFMGAN delivers more than 50\% RMSE reduction over
baselines across all landform types, with superior performance maintained across
spatial scales from 0.25 to 6\,m/pixel. The albedo-aware mechanism is directly
transferable to the dissertation problem: any generative method that synthesises
perspective-view surface appearance from nadir orbital data must similarly
disentangle shading caused by geometry from shading caused by intrinsic surface
properties. PFMGAN demonstrates this disentanglement is learnable from data
with appropriate physical constraints, providing a strong architectural prior.

Cao \textit{et al.}~\cite{cao2024gen3d} adopt a conditional GAN that maps
HiRISE orthorectified images to dense DEMs trained on paired HiRISE
orthoimage--DEM data at 2\,m and 0.25\,m resolution. Evaluation over four
geographically distinct Mars regions demonstrates height consistency between
predicted DEMs and stereo-derived HiRISE ground truth. The approach is closest
in spirit to the dissertation requirement: nadir imagery as input, a 3D
representation as output. However, the output is an elevation map, not a
perspective-rendered visual texture, and the evaluation geography is limited to
four areas with available stereo-based validation data.

Together, Cluster~4 establishes that the geometric component of the problem 
inferring surface shape from nadir images is tractable with physics-informed
deep learning. The component that remains unaddressed is rendering: given a
recovered surface geometry, synthesising what that surface would look like from
a rover's perspective, with correct photometric properties, realistic regolith
texture, and physically plausible shadows.

\subsection{Neural Scene Rendering for Martian Environments}
\label{subsec:nerf}

Neural Radiance Fields (NeRF)~\cite{paul2025nerfspace} represent a departure
from image-space processing: rather than enhancing or translating 2D images, NeRF
methods optimise a continuous volumetric scene function from which any viewpoint
can be rendered. This makes NeRF architecturally capable of the view-angle
transformation that SR and I2I methods cannot perform.

Giusti \textit{et al.}~\cite{giusti2022marf} demonstrate this capability in
MaRF, which applies NeRF to multi-image collections from the Curiosity rover at
Gale Crater. MaRF confirms that NeRF training converges on Martian imagery
despite the unusual photometric environment (persistent dust haze, low solar
elevation, high iron oxide spectral contrast), producing novel view synthesis
with PSNR and SSIM scores competitive with standard NeRF baselines on Earth
scenes. The work establishes Mars rover imagery as a viable NeRF training
domain. Its operational limitation is fundamental: MaRF requires multiple
overlapping rover images of the \textit{same surface patch} as input. For the
ExoMars RFM problem, real images of Oxia Planum do not yet exist, they will be
acquired after deployment. Any pre-mission VR training environment must
synthesise the rover-perspective view, not reconstruct it from data the rover
has not yet taken.

Li \textit{et al.}~\cite{li2025martian} extend neural Mars rendering into a
complete simulation pipeline through the Martian World Model system, comprising
two components. The M3arsSynth data engine processes NASA stereo navigation
imagery from Curiosity into a multimodal dataset of over 10,000 video clips
annotated with depth maps, surface normal estimates, and text descriptions.
The MarsGen video diffusion model, trained on this dataset, generates
controllable Martian landscape videos conditioned on text prompts, navigation
maps, and depth inputs. Experimental results demonstrate that MarsGen surpasses
video synthesis models trained on terrestrial datasets in terms of visual fidelity
and 3D structural consistency, confirming that domain-specific Mars training data
provides significant advantage over transfer from Earth imagery. MarsGen is the
most operationally complete system reviewed: it targets mission rehearsal and
robotic simulation as explicit use cases, and its controllable generation
addresses the requirement that operators can explore different traverse paths
in VR. Its limitation is the same as MaRF's, but at larger scale: the entire
training corpus is derived from rover perspective imagery. Terrain not previously
visited by any rover which encompasses virtually all of Mars outside four
landing sites, lies entirely outside MarsGen's generative support.

De Almeida \textit{et al.}~\cite{dealmeida2025neural3d} address a
geometrically adjacent problem: 3D reconstruction from descent-phase wide-angle
imagery. The descent trajectory provides a sequence of images spanning nadir to
oblique angles, but with strong radial distortion and limited baseline. The
paper demonstrates that an implicit neural height field representation
outperforms classical multi-view stereo (MVS) under these near-nadir, low-parallax
conditions, as the continuous surface prior suppresses depth estimation noise
that MVS accumulates. Evaluated on simulated descent sequences over high-fidelity
lunar and Mars terrain models, the neural approach achieves better RMSE across
a range of terrain roughness levels. This result matters for the dissertation in
a specific way: it demonstrates that neural representations with strong surface
priors can perform competitively in near-nadir imaging regimes where classical
geometric methods fail. If such representations could be coupled to a texture
synthesis module, the path from nadir orbital imagery to a rendered perspective
scene might run through an intermediate neural height-field representation
rather than through explicit 3D mesh construction.

Mart{\'i}nez-Petersen \textit{et al.}~\cite{martinezpetersen2025hifi} integrate
NeRF and Gaussian Splatting within a ROS~2-compatible pipeline that processes
raw rover rosbag recordings into dense photorealistic 3D reconstructions. The
key engineering contribution is Gaussian Splatting's render-time advantage: once
trained, a Gaussian scene representation renders at interactive frame rates
sufficient for real-time VR, unlike NeRF's ray-marching inference which requires
seconds per frame. The system requires perspective-view rover images as its
input; no mechanism is proposed or evaluated for deriving the Gaussian scene
from orbital data. The finding that Gaussian Splatting achieves VR-compatible
render rates is directly relevant to the dissertation's output target, however,
as it defines the representation format that the proposed nadir-to-perspective
synthesis pipeline should produce.

Paul \textit{et al.}~\cite{paul2025nerfspace} provide a comprehensive survey of
NeRF in space applications, identifying terrain mapping, navigation path planning,
and mission simulation as the three primary use cases. The survey notes that
existing NeRF methods in space contexts uniformly require multi-view imagery as
input either rover camera collections or descent image sequences and that the
problem of synthesising a NeRF from a single nadir orbital image has not been
addressed. This negative finding, reached by the most thorough treatment of the
topic, affirms the novelty of the proposed research direction.

\subsection{Virtual Reality Reconstruction and Mission Simulation}
\label{subsec:vr}

Catalano and Soccini~\cite{catalano2025inpainting} address a practical obstacle
in Martian VR construction: the heightmaps derived from CTX stereo and MOLA
altimetry contain systematic voids arising from instrument shadow, data dropout
during transmission, and incomplete orbital coverage. The proposed diffusion
model is trained on Mars terrain data and conditioned on the boundary of each
void region to inpaint missing elevation values in a structurally consistent
manner. Evaluation using SSIM and MS-SSIM against held-out regions of known
terrain demonstrates performance gains of up to 15\% in RMSE and 81\% in
perceptual LPIPS score over classical interpolation and prior deep learning
baselines. The approach genuinely improves heightmap completeness for VR terrain
meshes.

What the method does not produce is surface colour texture. A completed heightmap
defines the shape of the terrain mesh but assigns no RGB value to its vertices.
A VR environment built from inpainted heightmaps alone renders as a smooth,
textureless 3D surface: geometrically faithful but visually sterile and
photometrically unrepresentative of what a rover operator would actually see.
Texture synthesis, the step that converts a geometric skeleton into a
photorealistic surface is acknowledged but not performed. The paper explicitly
positions itself as addressing the geometry sub-problem, leaving the texture
synthesis challenge open.

Li \textit{et al.}~\cite{li2025martian} offer the architecturally most complete
VR simulation pipeline, but its dependence on rover perspective imagery as its
data substrate limits applicability to the four small geographic areas where
rover navigation cameras have operated. Pre-deployment simulation for Oxia
Planum, a site no rover has visited, lies outside its capability regardless of
generation quality.

\subsection{Hallucination and Epistemic Reliability of Generative Models}
\label{subsec:hallucination}

Every generative model reviewed in the preceding sections can, in principle,
produce plausible-looking surface features that have no basis in the underlying
physical reality. In natural image applications, invented texture is a quality
attribute. In mission-critical planetary science, a synthesised boulder that does
not exist or the absence of a real boulder that was not hallucinated which has
direct consequences for traversability assessment and crew safety.

Aithal \textit{et al.}~\cite{aithal2024hallucination} provide the most
systematic empirical analysis of this failure mode, identifying \textit{mode
interpolation} as its proximate mechanism. Diffusion models approximate the score
function of the training data distribution with a smooth neural network. When the
training distribution contains two distinct modes in close proximity for instance,
two crater morphologies that differ in ejecta structure. The smooth approximation
assigns non-zero probability to interpolated samples between those modes, which
never appeared in training data. Through controlled experiments on synthetic 2D
distributions and artificial image datasets, the paper shows that these
interpolated samples constitute hallucinations: the model synthesises features
that are statistically plausible under its learned score function but have no
physical referent in reality. Experiments on real-world image datasets confirm
the phenomenon at scale. In the Mars texture synthesis context, this means a
diffusion model trained on HiRISE imagery might synthesise impact crater rims
with morphologies that do not correspond to any crater at the target site,
aeolian ripple wavelengths inconsistent with the local wind regime, or rock
textures blending properties of geologically distinct surface units. The scale
of this risk is not bounded by model quality alone; it is a structural property
of smooth function approximation on multi-modal distributions.

Rathkopf~\cite{rathkopf2025reliability} reframes the hallucination problem as
an epistemological question: under what conditions can scientists rely on
generative model outputs for scientific inference? His answer draws on case
studies of AlphaFold~3 (protein structure prediction) and GenCast (atmospheric
modelling). In both cases, hallucination risk is converted to bounded,
quantifiable risk through two mechanisms: theory-guided training, where domain
physical laws impose constraints that prevent the model from generating
physically impossible outputs; and confidence-based error screening, where
uncertainty estimates flag outputs that should not be used for downstream
decisions. Rathkopf argues that generative models embedded in such structured
workflows support reliable scientific inference despite their opacity, provided
they operate in theoretically mature domains.

The implication for Mars texture synthesis is a concrete design requirement.
A generation pipeline that incorporates physical priors Mars solar irradiance
models, photometric functions for basaltic regolith, crater morphology
distributions derived from impact physics, and aeolian depositional
constraints cannot hallucinate a surface feature that violates those priors,
because the physics-violation incurs training loss. The dissertation must
therefore treat physical constraint not as an optional enhancement but as a
primary mechanism for bounding hallucination risk to scientifically acceptable
levels.

% =============================================================================
\section{Quantitative Synthesis and Comparative Analysis}
\label{sec:synthesis}

Table~\ref{tab:comparison} consolidates the twenty reviewed publications across
seven comparative dimensions: publication year, methodological approach, task
category, input-to-output view transformation, key reported performance metric,
whether the method was evaluated on Mars-specific data, and the core limitation
that prevents direct application to the nadir-to-rover-perspective VR synthesis
problem.

\begin{table*}[!t]
\caption{Comparative Overview of Reviewed Literature Against the Nadir-to-Rover-Perspective VR Synthesis Problem.
  Task codes: SR = super-resolution; I2I = image-to-image translation;
  DTM = digital terrain model generation; NeRF/GS = neural radiance field / Gaussian splatting;
  VR = VR reconstruction pipeline; Data = dataset contribution; Theory = theoretical analysis;
  Review = survey. View codes: N = nadir/orbital; P = perspective/rover; 3D = 3D representation;
  VR\textsubscript{H} = VR heightmap; ``Any'' = domain-agnostic.}
\label{tab:comparison}
\centering
\footnotesize
\setlength{\tabcolsep}{4pt}
\renewcommand{\arraystretch}{1.15}
\begin{tabular}{@{}p{0.45cm}lp{2.0cm}p{0.85cm}p{1.15cm}p{3.4cm}lp{4.4cm}@{}}
\toprule
\textbf{Ref.} & \textbf{Year} & \textbf{Method} & \textbf{Task} &
\textbf{View (I$\to$O)} & \textbf{Key Result} & \textbf{Mars?} &
\textbf{Core Limitation for VR Synthesis} \\
\midrule

\cite{tao2021cassis}    & 2021 & MARSGAN (AW-RRDB)             & SR      & N$\to$N   & ${\approx}3\times$ SR on CaSSIS; 8 sites, qualitative          & Yes     & No quantitative metric; nadir output only; ${\approx}1.3$\,m/pixel residual gap \\
\cite{wu2024rstsrn}     & 2024 & RSTSRN                        & SR      & N$\to$N   & ${\times}4$ SR; $+$0.88\,dB PSNR vs.\ LapSRN (${\times}4$)   & Yes     & Evaluated on bicubic degradation; real sensor PSF not modelled \\
\cite{lu2025thermal}    & 2025 & Multi-source Thermal SR        & SR      & N$\to$N   & ${\uparrow}$PSNR on THEMIS thermal; multi-modal fusion          & Yes     & Thermal domain only; no visual RGB texture output \\
\cite{khanna2024diffusionsat} & 2024 & DiffusionSat             & SR/Gen  & N$\to$N   & SOTA FID on multi-sensor RS; ICLR 2024                          & No      & No Mars training data; Earth geolocation conditioning \\
\cite{zhu2025taming}    & 2025 & TamingDiff-SRSISR             & SR      & N$\to$N   & ${\times}4$; ${\uparrow}$PSNR $+$ competitive FID on DOTA      & No      & Not Mars-specific; no view-angle transformation \\
\midrule
\cite{kim2024unsb}      & 2024 & UNSB (Schr{\"o}dinger Bridge) & I2I     & Any$\to$Any & ${\downarrow}$FID vs.\ CycleGAN/CUT; unpaired; ICLR 2024      & No      & Tested on colour-domain only; no geometry-aware view transform \\
\cite{xiao2025brownian} & 2025 & Dual-approx Bridge            & I2I     & Any$\to$Any & Deterministic I2I; ${\uparrow}$SSIM/LPIPS vs.\ stochastic; CVPR 2025 & No & No RS or planetary application demonstrated \\
\midrule
\cite{tao2021madnet}    & 2021 & MADNet 2.0                    & DTM     & N$\to$3D  & Monocular DTM ${\approx}$stereo accuracy; 25\,cm/pixel          & Yes     & Elevation only; no visual surface texture synthesis \\
\cite{zou2026pfmgan}    & 2026 & PFMGAN (physics-fusion GAN)   & DTM     & N$\to$3D  & $>{50\%}$ RMSE${\downarrow}$ over baselines; 5 landform types; 0.25--6\,m/pixel & Yes & DTM only; no perspective appearance rendering \\
\cite{cao2024gen3d}     & 2024 & Cond.\ GAN-DTM (ISPRS)        & DTM     & N$\to$3D  & Height consistency on 4 Mars regions vs.\ stereo ground truth   & Yes     & DEM output only; limited geographic validation \\
\cite{osadnik2025mcted} & 2025 & MCTED dataset                 & Data    & N (pairs) & 80,898 CTX-DEM pairs; ML-ready                                  & Yes     & Nadir$+$elevation only; no perspective-view samples \\
\cite{serrano2013hirise}& 2013 & Bayesian rock density (CTX)   & Analysis& N$\to$N   & Probabilistic rock density from CTX; Bayesian network           & Yes     & Density inference only; no texture or view transform \\
\midrule
\cite{giusti2022marf}   & 2022 & MaRF (NeRF)                   & NeRF    & P$\to$3D  & Novel view synthesis; PSNR/SSIM${\approx}$standard NeRF         & Yes     & Requires existing rover images; no synthesis from orbital data \\
\cite{li2025martian}    & 2025 & MarsGen / M3arsSynth          & NeRF+VR & P$\to$VR  & ${\uparrow}$FID vs.\ terrestrial video models; 10K+ video clips & Yes     & Rover perspective data only; unvisited terrain not covered \\
\cite{dealmeida2025neural3d} & 2025 & Neural height-field (descent) & NeRF & N/oblique$\to$3D & Neural HF$>$MVS on RMSE (simulated descent)               & Partial & Simulated data; descent-phase geometry only; no texture \\
\cite{martinezpetersen2025hifi} & 2025 & NeRF$+$GS pipeline (ROS~2) & NeRF+GS & P$\to$3D & Real-time GS render; ROS~2 integration; qualitative             & Partial & Rover perspective input only; no orbital-to-3D path \\
\cite{paul2025nerfspace}& 2025 & NeRF in space (survey)        & Review  & ---       & Comprehensive review; confirms no nadir$\to$perspective method  & Yes     & Survey only; confirms gap is open \\
\midrule
\cite{catalano2025inpainting} & 2025 & Diffusion inpainting (VR) & VR     & N$\to$VR\textsubscript{H} & ${\uparrow}$15\% RMSE; 81\% LPIPS vs.\ baselines; SSIM/MS-SSIM & Yes & Heightmap only; no RGB texture synthesis \\
\midrule
\cite{aithal2024hallucination} & 2024 & Mode interpolation analysis & Theory & --- & Mechanistic explanation of hallucination; NeurIPS 2024         & No      & Analytical only; no mitigation method proposed \\
\cite{rathkopf2025reliability} & 2025 & Hallucination-to-reliability & Theory & --- & Physics-guided training converts hallucination to bounded risk  & No      & Conceptual framework; no implementation or evaluation \\
\bottomrule
\end{tabular}
\end{table*}

Several patterns emerge from Table~\ref{tab:comparison} that are not apparent
from reading the papers individually. First, a clean partition separates
nadir-domain methods (all five SR papers and three DTM papers) from
perspective-domain methods (all three NeRF papers that were evaluated on real
Mars data). Not a single method crosses this boundary. The view angle of the
input determines the view angle of the output in every case except the two
theoretical I2I papers, which were benchmarked on natural images only.

Second, the methods with the strongest quantitative evidence of improvement
are those that incorporate explicit physical constraints. PFMGAN~\cite{zou2026pfmgan},
which embeds solar geometry and albedo modelling into its GAN, achieves more
than 50\% RMSE reduction over unconstrained baselines, the largest margin of any
generation method reviewed. By contrast, the SR papers relying on perceptual loss
functions without domain physics (MARSGAN, the early DiffusionSat applications)
produce plausible outputs but show greater tendency towards mode-interpolation
artefacts of the kind analysed in~\cite{aithal2024hallucination}. This is not a
coincidence: physical constraints restrict the generative model's output
distribution to a subspace consistent with planetary surface physics, and
hallucinated features that lie outside that subspace incur training loss.

Third, the two papers that come closest to a complete VR pipeline MarsGen
\cite{li2025martian} and the inpainting
system~\cite{catalano2025inpainting} each address one dimension of the
synthesis problem. MarsGen produces photorealistic perspective-view video but
requires rover imagery as input. The inpainting system completes elevation maps
from orbital data but produces no surface colour. A system that combines
orbital nadir input with perspective-view RGB texture output and is constrained
by physical priors does not exist in the reviewed literature.

Fourth, the two I2I papers~\cite{kim2024unsb, xiao2025brownian} represent the
strongest architectural candidates for closing the nadir-to-perspective gap,
but neither has been evaluated on remote sensing data. The jump from colour-domain
translation on natural images to geometry-dependent view-angle transformation on
satellite and rover imagery is substantial and constitutes an open research
problem.

% =============================================================================
\section{Open Challenges and Future Research Directions}
\label{sec:challenges}

\subsection{The Absence of Aligned Cross-View Training Data}
\label{subsec:data-gap}

Every supervised learning method reviewed requires either paired training samples
(nadir image and corresponding high-resolution target) or at minimum two separate
corpora that can be aligned statistically for unpaired training. For the
nadir-to-perspective Mars problem, no dataset provides geographically matched
nadir and perspective images of the same surface patch at aligned resolution.
MCTED~\cite{osadnik2025mcted} provides the largest nadir-CTX collection to date,
but paired with elevation maps rather than perspective imagery. The rover
navigation datasets used by MarsGen~\cite{li2025martian} provide perspective
views at four visited sites, with no co-registered nadir counterpart. Assembling
a training corpus that bridges both views whether through simulation, multi-view
geometric projection, or data augmentation is a methodological prerequisite
that none of the reviewed work addresses.

One route is synthetic data: projecting known Mars surface models (derived from
HiRISE stereo or from MADNet~2.0 monocular DTMs) into simulated rover
perspective renders. This approach is constrained by the realism of the Mars
texture models used to colour the geometry and by the domain gap between rendered
synthetic scenes and real rover imagery. Narrow-band domain gaps between
synthetic and real Mars images have been observed in landing hazard detection
contexts~\cite{serrano2013hirise}, and the gap between rendered and real surface
photometry is likely wider than domain adaptation methods such as
UNSB~\cite{kim2024unsb} have been tested on.

\subsection{Hallucination as a Safety-Critical Constraint}
\label{subsec:hallucination-challenge}

Mode interpolation~\cite{aithal2024hallucination} is a structural property of
diffusion-based generation, not a defect correctable by training longer or with
more data. For Earth-facing applications the consequences of a hallucinated
texture are aesthetic. For Mars mission training, a generated crater that does
not exist at the landing site, or the absence of a real hazard that the model
did not synthesise, could misdirect a crew's traverse planning. The reliability
framework of Rathkopf~\cite{rathkopf2025reliability} physics-guided training
plus confidence-based output screening provides the most credible route to
bounded hallucination risk, but it has not been instantiated for planetary
imagery generation. Key domain priors that would constrain Mars generation
include: cratering density and size-frequency distributions from impact physics;
aeolian bedform wavelength--height relationships; photometric functions for basaltic
and ferric oxide surface compositions; and shadow angle constraints imposed by
Mars's axial tilt and the landing site's latitude and season. Embedding these in
the training objective is an open engineering problem.

\subsection{The Computational Cost of Diffusion Synthesis at Planetary Scale}
\label{subsec:compute}

Diffusion models require iterative denoising over many steps (typically 20--1000)
per generated sample. DiffusionSat~\cite{khanna2024diffusionsat} operates on
individual tiles at fixed resolution; MarsGen~\cite{li2025martian} generates
video clips of limited spatial extent. Generating continuous,
full-site Mars surface texture at rover camera resolution would require either
tiling with consistent boundary conditions across tiles, or a cascaded coarse-to-fine
approach. Neither strategy has been demonstrated at the scale needed for a
mission landing site: Oxia Planum spans roughly $200 \times 200$\,km, and even
a modest 100-pixel-per-metre rover texture resolution implies generating $4 \times
10^{10}$ RGB pixels. Memory requirements for a diffusion model trained at this
scale, and inference times for generating the full canvas, represent unsolved
engineering problems. The Gaussian Splatting representation demonstrated
in~\cite{martinezpetersen2025hifi} offers one path toward efficient storage and
real-time rendering once the scene has been generated, but it does not address
the generation bottleneck.

\subsection{Evaluation Benchmarks for Mars Surface Generation}
\label{subsec:eval}

The reviewed SR literature uses at least four distinct evaluation protocols:
qualitative geomorphological expert assessment~\cite{tao2021cassis}, PSNR/SSIM
against bicubically downsampled originals~\cite{wu2024rstsrn}, perceptual metrics
(LPIPS, FID) on remote sensing benchmarks~\cite{zhu2025taming, khanna2024diffusionsat},
and structural metrics (SSIM, MS-SSIM) on heightmap inpainting~\cite{catalano2025inpainting}.
No standardised benchmark exists for evaluating the photorealistic quality and
physical correctness of synthesised Mars perspective-view imagery. Without such
a benchmark, comparing methods or demonstrating progress towards the VR synthesis
objective is not possible in a reproducible manner. Establishing a Mars texture
synthesis benchmark with held-out rover images as reference targets and a
protocol for generating nadir inputs at matched locations from orbital
archives is a necessary community infrastructure contribution.

\subsection{Multi-Modal Fusion Across Instrument Types}
\label{subsec:multimodal}

The thermal SR work of Lu and Su~\cite{lu2025thermal} and the Bayesian rock
density prediction of Serrano \textit{et al.}~\cite{serrano2013hirise} both
demonstrate that information from multiple Mars instruments can be combined to
recover surface properties unresolvable from any single data source. Visible
imaging, thermal IR, and topographic data each carry distinct information
about surface composition, thermophysical properties, and physical structure. A
generative pipeline that fuses all three using thermal emission to constrain
surface material type, topographic derivatives to constrain shape-from-shading
ambiguities, and visible reflectance to anchor albedo which would face far fewer
hallucination degrees of freedom than a model trained on visible imagery alone.
No reviewed paper pursues this multi-modal fusion strategy for texture synthesis,
and MCTED's limitation to visible-elevation pairs underscores the data gap this
would require addressing.

% =============================================================================
\section{Conclusion}
\label{sec:conclusion}

Twenty publications reviewed here collectively demonstrate that the constituent
capabilities needed for VR-grade Mars landing site visualisation, super-resolution
of orbital imagery, image domain translation, DTM generation, neural scene
rendering, and diffusion-based inpainting have each advanced substantially
since 2021. Transformer-based and diffusion-based SR methods now reliably
outperform GAN baselines on orbital Mars imagery. Physics-constrained GANs
recover Martian terrain geometry with more than 50\% lower reconstruction error
than their unconstrained predecessors. NeRF and Gaussian Splatting have been
demonstrated on real Mars rover data and can, in principle, produce real-time
renderable scenes. Diffusion inpainting completes missing elevation maps with
quantifiably superior structural fidelity.

What has not been demonstrated is the integration of these capabilities into
a pipeline that starts from nadir orbital data and ends with rover-perspective
visual texture at VR-grade fidelity. The view-angle transformation from
nadir to perspective is the missing operation, and it appears in none of the
reviewed methods. The I2I frameworks of~\cite{kim2024unsb} and~\cite{xiao2025brownian}
are architecturally the strongest candidates for this transformation, but neither
has been applied to remote sensing imagery, let alone to the photometrically
complex domain shift between nadir Mars orbital and perspective rover views.

The proposed dissertation targets this specific gap. The analysis of
hallucination risk in~\cite{aithal2024hallucination} and the reliability
framework of~\cite{rathkopf2025reliability} together establish that
physics-constrained training is not optional but mandatory for the outputs to
be scientifically usable in mission planning. PFMGAN's result~\cite{zou2026pfmgan}
that explicit physical priors reduce reconstruction error by more than half which
provides a proof-of-principle that such constraints are learnable from data.
The combination of unpaired I2I translation (to bridge the view domains without
requiring aligned training pairs), physics-informed training objectives (to bound
hallucination risk), and a Gaussian Splatting output representation (to enable
real-time VR rendering) defines the most promising direction identified by this
survey, and it is the direction the proposed work will pursue.

% =============================================================================
\bibliographystyle{IEEEtran}
\bibliography{refs}

\end{document}
